\subsection{17.03.2026}

\subsubsection{Ziel}

Ziel des Tages war es, den zweiten DHT11-Sensor
anzuschließen und in die Taupunktberechnung zu integrieren.

\subsubsection{Durchgeführte Arbeiten}

Zunächst wurde der zweite DHT11-Sensor an Pin 37
(GPIO 26) angeschlossen und die Verkabelung überprüft.

Anschließend wurde in der Datei
\texttt{dht11\_jh.py} die Klasse
\texttt{dht11Sensor} erweitert, sodass die
GPIO-Pin-Nummer beim Erstellen eines Objekts
übergeben werden kann.

Danach wurde in
\texttt{taupunkt\_main.py} ein zweiter
Sensor mit der entsprechenden GPIO-Nummer
initialisiert.

Zusätzlich wurde eine Konsolenausgabe ergänzt,
welche bei ungültigen oder noch nicht verfügbaren
Messwerten die Meldung „messung...“
ausgibt.

Zum Ende des Tages wurde die bestehende
Taupunktberechnung in die Hauptlogik integriert.

\subsubsection{Problem 1: Falsche GPIO-Pin-Nummer}

Zunächst wurde für den zweiten Sensor die
physische Pin-Nummer 37 direkt verwendet.

Dies führte dazu, dass sich die Verbindung
über VNC aufhing und der Raspberry Pi erst nach
einem Neustart wieder erreichbar war.

\textbf{Lösung}

Es wurde festgestellt, dass im verwendeten
GPIO-Modus nicht die physische Pin-Nummer,
sondern die GPIO-Nummer verwendet werden muss.

Da im Projekt der BCM-Modus genutzt wird,
musste für den Sensor die GPIO-Nummer 26
anstatt der physischen Pin-Nummer 37 verwendet
werden.

\subsubsection{Problem 2: Fehler in der Taupunktberechnung}

Beim Ausführen der Taupunktberechnung trat die
Fehlermeldung
\texttt{ValueError: math domain error}auf.

Der Fehler entstand bei der Berechnung mit
\texttt{log10(...)}.

\textbf{Lösung}

Es wurde festgestellt, dass ungültige oder
unvollständige Messwerte an die Berechnung
übergeben wurden.

Dadurch entstand ein mathematisch nicht
zulässiger Wert innerhalb der Logarithmusfunktion.

Die Ursache soll in der nächsten
Unterrichtseinheit weiter untersucht und
behoben werden.

\subsubsection{Erkenntnisse}

Es wurde festgestellt, dass bei Verwendung des
BCM-Modus ausschließlich GPIO-Nummern und
nicht die physischen Pin-Nummern verwendet
werden dürfen.

Außerdem wurde deutlich, dass Sensordaten vor
der Weiterverarbeitung auf gültige Werte
überprüft werden müssen, um Fehler bei
mathematischen Berechnungen zu vermeiden.

\subsubsection{Nächste Schritte}

Als nächster Schritt soll die Taupunktberechnung
fehlerfrei umgesetzt werden.
